\chapter{Objetivos}
\label{ch:objetivos}

Los objetivos que se enumeran a continuación proceden de la propuesta de Trabajo
Fin de Grado aprobada y se reproducen sin modificaciones de redacción. El
numerado automático de la lista sustituye a la numeración manual del documento
original.

\begin{enumerate}[label=\arabic*.,leftmargin=*]
  \item Diseñar e implementar un dashboard interactivo que integre datos de
        rendimiento obtenidos de plataformas y dispositivos deportivos (como
        Strava o Garmin), permitiendo la visualización dinámica de métricas
        individuales y grupales.

  \item Desarrollar un modelo de análisis de datos deportivos que permita
        detectar patrones de evolución en el ritmo, frecuencia cardíaca,
        potencia y distancia recorrida, utilizando técnicas estadísticas y de
        ingeniería de datos.

  \item Implementar un módulo predictivo basado en regresión lineal y modelos de
        series temporales para estimar la progresión del rendimiento del
        corredor a corto y medio plazo.

  \item Diseñar una interfaz visual optimizada, intuitiva y adaptable, que
        facilite la interpretación de los indicadores de rendimiento tanto para
        deportistas individuales como para entrenadores o equipos.

  \item Integrar un sistema de gestión de datos estructurados que permita la
        comparación de resultados entre distintos corredores, fomentando la toma
        de decisiones basada en evidencias y el seguimiento personalizado del
        progreso.
\end{enumerate}

\pendiente{cada objetivo deberá trazarse posteriormente a los requisitos del
capítulo de análisis del problema y a la sección de verificación
correspondiente, de modo que sea posible comprobar su cumplimiento uno a uno.}
